\section{Conclusão e trabalhos
futuros}\label{conclusuxe3o-e-trabalhos-futuros}

Este trabalho desenvolveu um pipeline reprodutível para investigar a
intensidade de reutilização de Agent Skills como artefatos de engenharia
de software. A pesquisa modelou o número de repositórios que contêm
exatamente o mesmo \texttt{SKILL.md}, condicionando a análise às skills
reutilizadas. O protocolo combinou atributos estruturais, sinais
textuais e contexto de repositório, utilizou validação agrupada por
\texttt{representative\_repo\_id} e realizou a seleção de atributos
dentro das dobras.

Seis regressores foram comparados sobre 25.918 skills reutilizadas na
visão \texttt{repository-only}. O Random Forest foi escolhido pelo menor
RMSLE médio na validação agrupada (0,5753) e obteve 0,5833 no teste
externo de 6.479 skills. O Gradient Boosting teve MAE externo menor
(3,1820 contra 3,2663). Uma referência constante baseada na mediana do
desenvolvimento teve RMSLE pior (0,7142), mas MAE melhor (3,0665) que o
Random Forest. A escolha do modelo depende, portanto, do objetivo e da
métrica.

A previsão da reutilização é parcial. O Random Forest subestimou todas
as 324 skills do teste acima do P95 do desenvolvimento e teve MAE de
35,15 nessa faixa. O Ridge híbrido apresentou falha semelhante na
validação interna, em análise separada. A comparação das três
representações com Ridge sugere sinal em texto e estrutura, mas a
diferença entre \texttt{hybrid} e \texttt{skill-only} é pequena diante
da variação entre dobras e ainda não foi testada externamente. As
importâncias por permutação do Random Forest se referem a contexto de
repositório, não a propriedades intrínsecas da skill.

Concluímos que características observáveis permitem prever parte da
intensidade de replicação exata entre repositórios, sobretudo na região
central. A evidência não permite explicar por que uma skill se
dissemina, inferir qualidade ou formular uma regra para torná-la mais
reutilizável.

Estudos futuros podem avaliar as representações textuais em um novo
teste externo ainda não consultado, comparar outras estratégias para a
cauda e validar hipóteses em outros snapshots ou ecossistemas. O
presente relatório documenta os resultados e os limites do protocolo
executado.
